% !TEX program = lualatex
\documentclass[12pt, a4paper]{article}
\usepackage{master}
\usepackage{tikz}
\usetikzlibrary{calc, positioning, arrows.meta}

\DeclareWorkGR{Cat}{범주론}{Κατηγορίαι}{Categoriae}
\DeclareWorkGR{Top}{토피카}{Τοπικά}{Topica}
\DeclareWorkGR{Met}{형이상학}{Μεταφυσικά}{Metaphysica}
\DeclareWorkGR{APo}{분석론 후서}{Ἀναλυτικὰ ὕστερα}{Analytica posteriora}

\fancyhead[L]{\footnotesize\color{midgray}오르가논 강독}
\fancyhead[R]{\footnotesize\color{midgray}제6주}

\begin{document}
\thispagestyle{firstpage}

\weeksection{제6주}{이사고게 I--III장}{서론, 류, 종}

\subsection{이사고게란 무엇인가}

\subsubsection{저자와 배경}

이번 주부터 우리는 아리스토텔레스의 텍스트를 잠시 떠나, 포르피리오스(\gr{Πορφύριος}, c.\ 234--305)의 이사고게(\gr{Εἰσαγωγή}, `입문')를 읽습니다. 이사고게는 약 268--270년에 시칠리아에서 크뤼사오리오스(\gr{Χρυσαόριος})라는 인물에게 보내는 형식으로 작성되었습니다. 분량은 부세판(Busse)으로 겨우 20쪽 남짓에 불과하지만, 이 짧은 텍스트는 서양 논리학사에서 가장 오래 살아남은 교과서가 됩니다 --- 6세기부터 16세기까지, 비잔틴·이슬람·라틴 서방의 세 문명권에서 철학 교육을 시작하는 모든 학생이 이 텍스트를 읽었습니다.

포르피리오스는 신플라톤주의(Neoplatonism)의 핵심 인물로, 플로티노스(Plotinos)의 제자이자 편집자입니다 --- 플로티노스의 엔네아데스(\gr{Ἐννεάδες})를 편집한 사람이 바로 포르피리오스입니다. 그러나 이사고게는 놀라울 정도로 신플라톤주의적 색채가 엷습니다. 에밀손(Emilsson)이 지적하듯이, 이것이 바로 이 텍스트가 서로 다른 철학적 전통에 걸쳐 채택될 수 있었던 비결이기도 합니다.\footnote{Eyjólfur Emilsson, ``Porphyry,'' in \gri{Stanford Encyclopedia of Philosophy}, rev.\ 2021. 에밀손은 이사고게에서 `독특하게 플라톤적인 견해의 부재'와 `보편자의 존재론적 지위에 관한 깊은 물음들의 회피'를 강조합니다.}

\subsubsection{이사고게의 목적}

이사고게의 첫 문장은 텍스트의 목적을 선언합니다.\footnote{이사고게 그리스어 텍스트: Adolf Busse, ed., \gri{Porphyrii Isagoge et in Aristotelis Categorias Commentarium}, CAG IV.1 (Berlin: G.\ Reimer, 1887). Eng.\ trans.: Jonathan Barnes, \gri{Porphyry: Introduction}, Clarendon Later Ancient Philosophers (Oxford: Clarendon Press, 2003). 본 강의록의 이사고게 인용은 부세판 쪽·행 번호를 따릅니다.}

\srcquote{\gr{Ὄντος ἀναγκαίου, Χρυσαόριε, καὶ εἰς τὴν τῶν παρὰ Ἀριστοτέλει κατηγοριῶν διδασκαλίαν τοῦ γνῶναι τί γένος καὶ τί διαφορὰ τί τε εἶδος καὶ τί ἴδιον καὶ τί συμβεβηκός \ldots σύντομόν σοι παράδοσιν ποιούμενος πειράσομαι διὰ βραχέων ὥσπερ ἐν εἰσαγωγῆς τρόπῳ τὰ παρὰ τοῖς πρεσβυτέροις ἐπελθεῖν, τῶν μὲν βαθυτέρων ἀπεχόμενος ζητημάτων \ldots}}{`크뤼사오리오스여, 아리스토텔레스의 범주론을 배우기 위해 류와 종차와 종과 고유성과 부수성이 무엇인지를 아는 것이 필요하므로 \ldots 나는 그대에게 입문의 방식으로 선배들이 말한 바를 짧게 살펴보려 한다. 더 깊은 탐구들은 삼가면서 \ldots'}{이사고게 1.1--8 Busse}

여기서 주목할 점이 네 가지 있습니다. 첫째, 포르피리오스는 이 텍스트의 목적을 범주론 학습의 예비 작업으로 명시합니다 --- 다만 반스(Barnes 2003)는 이사고게가 범주론만이 아니라 `논리학 일반'의 입문서라고 주장합니다. 정의(\gr{ὁρισμός}), 분할(\gr{διαίρεσις}), 논증(\gr{ἀπόδειξις})이 모두 언급되는데, 이것들은 범주론이 아니라 분석론 후서와 토피카의 주제이기 때문입니다. 둘째, 다섯 가지 술어가능자(\gr{πέντε φωναί}, `다섯 목소리')가 이미 첫 문장에 전부 열거됩니다. 셋째, `더 깊은 탐구들'을 의도적으로 회피하겠다는 선언이 나옵니다. 넷째, `선배들'(\gr{τοῖς πρεσβυτέροις})과 특히 `페리파토스학파'(\gr{οἱ ἐκ τοῦ περιπάτου})를 따르겠다는 선언이 텍스트의 아리스토텔레스적 지향을 표명합니다.

\subsubsection{보편자에 관한 세 가지 물음 --- 포르피리오스의 거부}

첫 문장 직후(1.9--16 Busse), 포르피리오스는 류와 종의 존재론적 지위에 관한 세 가지 물음을 제기합니다 --- 그리고 대답하기를 거부합니다.

\srcquote{\gr{αὐτίκα περὶ τῶν γενῶν τε καὶ εἰδῶν τὸ μὲν εἴτε ὑφέστηκεν εἴτε καὶ ἐν μόναις ψιλαῖς ἐπινοίαις κεῖται εἴτε καὶ ὑφεστηκότα σώματά ἐστιν ἢ ἀσώματα καὶ πότερον χωριστὰ ἢ ἐν τοῖς αἰσθητοῖς καὶ περὶ ταῦτα ὑφεστῶτα, παραιτήσομαι λέγειν βαθυτάτης οὔσης τῆς τοιαύτης πραγματείας καὶ ἄλλης μείζονος δεομένης ἐξετάσεως.}}{`류와 종에 관해 --- (a) 그것들이 존립하는지, 아니면 벌거벗은 사유에만 놓여 있는지; (b) 만약 존립한다면, 물체적인지 비물체적인지; (c) 감각적인 것들과 분리되어 있는지, 아니면 감각적인 것들 안에 그리고 그것들에 관하여 존립하는지 --- 나는 말하기를 사양하겠다. 이런 종류의 탐구는 가장 깊은 것이며, 다른, 더 큰 검토를 필요로 하기 때문이다.'}{이사고게 1.9--16 Busse}

이 세 물음은 조건적 계열을 이룹니다 --- 각 물음은 앞 물음에 긍정적으로 대답해야만 의미가 있습니다. 물음 (a)는 실재론과 유명론의 대립을 준비하고, 물음 (b)는 스토아학파(물체적)와 플라톤·아리스토텔레스(비물체적) 사이의 대립을 반영하며, 물음 (c)는 플라톤적 실재론(분리된 형상)과 아리스토텔레스적 실재론(감각적인 것들 안에 있는 형상)의 대립을 준비합니다. 이 물음들은 중세 `보편자 논쟁'(universals debate)의 직접적인 기원입니다.\footnote{SEP, ``The Medieval Problem of Universals'' (Gyula Klima, rev.\ 2022).}

포르피리오스가 대답을 거부한 이유에 대해 학자들의 의견이 갈립니다. 반스(Barnes 2003)와 에베센(Ebbesen 1990)은 포르피리오스가 논리학을 존재론적으로 중립적인 학문으로 간주했다고 봅니다. 반면 키아라돈나(Chiaradonna 2008)는 `논리학자 포르피리오스를 철학자 포르피리오스로부터 그렇게 쉽게 분리할 수 없다'고 반박합니다.\footnote{Riccardo Chiaradonna, ``What is Porphyry's Isagoge?,'' \gri{Documenti e studi sulla tradizione filosofica medievale} 19 (2008): 1--30.}

여기서 사고 연습을 해봅시다. `사람'이라는 종(\gr{εἶδος})은 소크라테스, 플라톤 등 개별 인간들에 대해 술어됩니다. 그렇다면 `사람'이라는 종은 소크라테스와 플라톤 외에도 무엇인가로서 존재하는 것일까요? 범주론에서 다루었던 `제2실체'의 지위 문제가 여기서 새로운 차원으로 제기됩니다. 범주론 5장에서 아리스토텔레스는 제2실체가 `주어에 대해 말해지는 것'이라고 했지만, 제2실체가 제1실체와 독립적으로 존재하는지의 물음은 열어 두었습니다. 포르피리오스의 세 물음은 바로 범주론이 열어 둔 이 빈자리를 정면으로 겨냥합니다.

\subsubsection{아리스토텔레스의 네 술어가능자와 포르피리오스의 다섯 술어가능자}

포르피리오스는 다섯 술어가능자의 예시를 제시합니다(2.15--20 Busse).

\srcquote{\gr{ἔστι δὲ γένος μὲν οἷον τὸ ζῷον, εἶδος δὲ οἷον ὁ ἄνθρωπος, διαφορὰ δὲ οἷον τὸ λογικόν, ἴδιον δὲ οἷον τὸ γελαστικόν, συμβεβηκὸς δὲ οἷον τὸ λευκόν, τὸ μέλαν, τὸ καθέζεσθαι.}}{`류는 예컨대 동물이고, 종은 예컨대 사람이며, 종차는 예컨대 이성적인 것이고, 고유성은 예컨대 웃을 수 있음이며, 부수성은 예컨대 흼, 검음, 앉아 있음이다.'}{이사고게 2.15--20 Busse}

아리스토텔레스의 \citework{Top}{I.4, 101b17--25}에는 네 가지 술어가능자가 있습니다 --- 정의(\gr{ὅρος}), 류(\gr{γένος}), 고유성(\gr{ἴδιον}), 부수성(\gr{συμβεβηκός}). 포르피리오스의 목록은 아리스토텔레스의 `정의'를 `류 + 종차 → 종'으로 분해한 것입니다. 이것은 논리적으로 정합적입니다 --- 아리스토텔레스적 논리학에서 정의는 류 + 종차이므로, 정의를 구성 요소로 분해하면 류, 종차, 종이 됩니다.\footnote{Christos Evangeliou, ``Aristotle's Doctrine of Predicables and Porphyry's \gri{Isagoge},'' \gri{Journal of the History of Philosophy} 23 (1985): 15--34.}

아리스토텔레스의 네 가지 분류 기준도 명확히 해둡시다. 아리스토텔레스는 두 가지 기준으로 분류합니다: 술어가 주어에 대해 본질적인가(\gr{ἐν τῷ τί ἐστι}) 여부와, 술어가 주어와 외연이 같은가(\gr{ἀντιστρέφει}) 여부입니다. 정의는 본질적이면서 외연이 같습니다. 류와 종차는 본질적이지만 외연이 같지 않습니다. 고유성은 본질적이지 않지만 외연이 같습니다. 부수성은 본질적이지도 않고 외연이 같지도 않습니다.

\begin{concepttable}{|c|c|c|c|c|}
\hline
\rowcolor{tableheadblue}
{\color{h1navy}\bfseries} &
{\color{h1navy}\bfseries 본질적?} &
{\color{h1navy}\bfseries 역전?} &
{\color{h1navy}\bfseries 아리스토텔레스} &
{\color{h1navy}\bfseries 포르피리오스} \\ \hline
류 & 예 & 아니오 & 류 & 류 \\ \hline
종차 & 예 & 아니오 & (정의의 일부) & 종차 \\ \hline
종 & 예 & — & (정의의 결과) & 종 \\ \hline
고유성 & 아니오 & 예 & 고유성 & 고유성 \\ \hline
부수성 & 아니오 & 아니오 & 부수성 & 부수성 \\ \hline
\end{concepttable}

\subsection{류 (II장)}

\subsubsection{류의 정의}

포르피리오스는 류를 다음과 같이 정의합니다(2.22--3.1 Busse).

\srcquote{\gr{γένος εἶναι λέγοντες τὸ κατὰ πλειόνων καὶ διαφερόντων τῷ εἴδει ἐν τῷ τί ἐστι κατηγορούμενον, οἷον τὸ ζῷον.}}{`류란 종에 있어서 서로 다른 여럿에 대해 그것이 무엇인지를 말해주며 술어되는 것을 말한다. 예컨대 동물.'}{이사고게 2.22--3.1 Busse}

핵심은 \gr{ἐν τῷ τί ἐστι}(`그것이 무엇인지를 말해주며')라는 조건입니다. `소크라테스는 무엇인가?'라고 물었을 때 `동물'이라고 대답할 수 있으며, 이때 `동물'은 류로서 기능합니다. 류는 양이나 질이 아니라 본질의 범주에서 술어됩니다 --- 이것은 범주론 5장에서 제2실체가 제1실체에 대해 `무엇인지'를 말해주는 것과 정확히 대응합니다.

또한 `여럿에 대해'(\gr{κατὰ πλειόνων})라는 조건도 중요합니다. 류는 반드시 복수의 종에 대해 술어되어야 합니다 --- `동물'은 사람, 말, 개 등 여러 종에 술어됩니다.

\subsubsection{류의 위계와 포르피리오스의 나무}

포르피리오스는 위계의 양 끝을 구별합니다. 최고류(\gr{γενικώτατον γένος}) --- 예컨대 실체(\gr{οὐσία}) --- 는 류이기만 하고 어떤 것의 종이 되지 않습니다. 최저종(\gr{εἰδικώτατον εἶδος}) --- 예컨대 사람 --- 은 종이기만 하고 더 이상 류가 되지 않습니다. 그 중간에 있는 하위류(\gr{ὑπάλληλα γένη})는 아래를 향해서는 류이고 위를 향해서는 종입니다.

\srcquote{\gr{ἡ οὐσία ἔστι μὲν καὶ αὐτὴ γένος, ὑπὸ δὲ ταύτην ἐστὶν σῶμα, καὶ ὑπὸ τὸ σῶμα ἔμψυχον σῶμα, ὑφ᾽ ὃ τὸ ζῷον, ὑπὸ δὲ τὸ ζῷον λογικὸν ζῷον, ὑφ᾽ ὃ ὁ ἄνθρωπος, ὑπὸ δὲ τὸν ἄνθρωπον Σωκράτης καὶ Πλάτων καὶ οἱ κατὰ μέρος ἄνθρωποι.}}{`실체는 그 자체로 류이다. 그 아래에 물체가 있고, 물체 아래에 생명 있는 물체가, 그 아래에 동물이, 동물 아래에 이성적 동물이, 그 아래에 사람이, 사람 아래에 소크라테스와 플라톤과 개별 사람들이 있다.'}{이사고게 4.21 이하 Busse}

이 대목이 `포르피리오스의 나무'(arbor Porphyriana)의 텍스트적 기초입니다.

% ── 포르피리오스의 나무 tikz (비플로트, 본문 인라인) ──
\par\nobreak\vspace{6pt}
\begin{center}
\begin{tikzpicture}[
  scale=0.72, every node/.style={transform shape},
  box/.style={draw, rounded corners=2pt, minimum width=34mm, minimum height=8mm, align=center, font=\footnotesize},
  centerbox/.style={draw, rounded corners=2pt, minimum width=34mm, minimum height=8mm, align=center, fill=tableheadblue, font=\footnotesize\bfseries},
  indiv/.style={draw, rounded corners=2pt, minimum width=34mm, minimum height=7mm, align=center, font=\scriptsize, dashed},
  arr/.style={-{Stealth[length=4pt]}, thick, color=midgray},
  lbl/.style={font=\scriptsize, color=midgray, midway}
]
  \node[centerbox] (sub) {{\color{h1navy}실체} (\gr{οὐσία})};
  \node[box, below=10mm of sub] (body) {물체 (\gr{σῶμα})};
  \node[box, below=10mm of body] (anim) {생명 있는 물체 (\gr{ἔμψυχον σῶμα})};
  \node[box, below=10mm of anim] (animal) {동물 (\gr{ζῷον})};
  \node[box, below=10mm of animal] (rational) {이성적 동물 (\gr{λογικὸν ζῷον})};
  \node[centerbox, below=10mm of rational] (human) {{\color{h1navy}사람} (\gr{ἄνθρωπος})};
  \node[indiv, below=10mm of human] (indivs) {소크라테스, 플라톤 \ldots};

  \draw[arr] (sub) -- node[right, lbl] {물체적/비물체적} (body);
  \draw[arr] (body) -- node[right, lbl] {생명 있는/없는} (anim);
  \draw[arr] (anim) -- node[right, lbl] {감각 있는/없는} (animal);
  \draw[arr] (animal) -- node[right, lbl] {이성적/비이성적} (rational);
  \draw[arr] (rational) -- (human);
  \draw[arr, dashed] (human) -- (indivs);

  % 왼쪽 레이블
  \node[left=18mm of sub, font=\scriptsize\itshape, color=midgray] {최고류};
  \node[left=18mm of human, font=\scriptsize\itshape, color=midgray] {최저종};
  \node[left=18mm of indivs, font=\scriptsize\itshape, color=midgray] {개별자};
\end{tikzpicture}

{\small\color{midgray} 포르피리오스의 나무: 실체에서 개별자까지의 류-종 위계. 원문에 도식은 없으며, 중세의 창작물이다.}
\end{center}
\vspace{-2pt}

여기서 중요한 학술적 사실을 짚어야 합니다: 이사고게 원문에는 나무 도식이 전혀 등장하지 않습니다. 고대 주석자들은 이 위계를 `사다리'(\gr{κλίμαξ})라고 불렀으며, 나무 도식은 보에티우스의 주석서 필사본에서 처음 등장한 중세의 창작물입니다. `포르피리오스의 나무'라는 이름은 13세기 페트루스 히스파누스(Petrus Hispanus)의 논리학 총론에서 확립되었습니다.\footnote{Annemieke Verboon, ``The Medieval Tree of Porphyry,'' in \gri{The Tree: Symbol, Allegory and Structural Device in Medieval Art and Thought}, International Medieval Research 20 (Turnhout: Brepols, 2014), 83--101.}

포르피리오스는 또한 아리스토텔레스(\citework{Met}{998b22--23})를 따라, 존재(\gr{τὸ ὄν}) 자체는 모든 것에 공통된 하나의 류가 아니라고 말합니다. 류가 종으로 분할되려면 종차가 필요한데, 종차는 류 자체에 참여하지 않아야 합니다. 그런데 모든 것이 `존재'에 참여하므로, 존재를 분할할 종차가 존재의 바깥에 있을 수 없습니다 --- 따라서 존재는 류가 아닙니다.

\subsection{종 (III장)}

\subsubsection{종의 두 가지 정의}

포르피리오스는 종에 대해 두 가지 정의를 제시합니다. 관계적 정의: `종은 류 아래에 배열되는 것이며, 류가 그것에 대해 무엇인지를 말해주며 술어되는 것이다.' 최저종의 정의:

\srcquote{\gr{εἶδός ἐστι τὸ κατὰ πλειόνων καὶ διαφερόντων τῷ ἀριθμῷ ἐν τῷ τί ἐστι κατηγορούμενον.}}{`종은 수에 있어서 서로 다른 여럿에 대해 그것이 무엇인지를 말해주며 술어되는 것이다.'}{이사고게 III장}

류와 종의 정의를 비교하면 결정적인 차이가 드러납니다: 류는 `종에 있어서 서로 다른'(\gr{τῷ εἴδει}) 여럿에 술어되고, 종은 `수에 있어서 서로 다른'(\gr{τῷ ἀριθμῷ}) 여럿에 술어됩니다. `동물'(류)은 사람, 말, 개 등 서로 종이 다른 여럿에 술어됩니다. `사람'(종)은 소크라테스, 플라톤 등 서로 수만 다른 여럿에 술어됩니다.

\subsubsection{종의 중간적 지위와 범주론과의 연결}

최저종은 본질적 술어화가 작동하는 가장 낮은 수준입니다. 이것은 2주차에서 다루었던 범주론의 제1실체·제2실체 구별과 직접 연결됩니다 --- 범주론 5장에서 아리스토텔레스는 `종이 류보다 더 실체적'이라고 했습니다(2b7--14). 포르피리오스의 최저종 개념은 이 직관을 체계화합니다 --- 최저종은 개별자에 가장 가까운 본질적 술어이며, 따라서 가장 `적합한' 답변을 제공합니다.

여기서 사고 연습을 해봅시다. 포르피리오스의 나무에서 각 단계는 종차에 의해 분할됩니다. 그런데 이 분할이 완전한지(exhaustive) 그리고 배타적인지(exclusive) 어떻게 보장할 수 있을까요? `동물'을 `이성적/비이성적'이 아니라 `육상/수중/공중'으로 분할할 수도 있지 않을까요? 아리스토텔레스 자신도 동물부분론과 동물지에서 다양한 분류 기준을 사용합니다. 포르피리오스의 나무는 하나의 분할 경로만을 보여주지만, 실제로는 같은 류를 여러 종차에 의해 서로 다른 방식으로 분할할 수 있다는 문제가 이미 여기에 잠재해 있습니다.

% ── 참고문헌 ──
\subsection*{참고문헌}
\begin{references}

\refitem Barnes, Jonathan. \gri{Porphyry: Introduction}. Clarendon Later Ancient Philosophers. Oxford: Clarendon Press, 2003.

\refitem Busse, Adolf, ed. \gri{Porphyrii Isagoge et in Aristotelis Categorias Commentarium}. CAG IV.1. Berlin: G.\ Reimer, 1887. Eng.\ trans.: Jonathan Barnes, \gri{Porphyry: Introduction}, Clarendon Later Ancient Philosophers. Oxford: Clarendon Press, 2003.

\refitem Chiaradonna, Riccardo. ``What is Porphyry's Isagoge?'' \gri{Documenti e studi sulla tradizione filosofica medievale} 19 (2008): 1--30.

\refitem Emilsson, Eyjólfur. ``Porphyry.'' In \gri{Stanford Encyclopedia of Philosophy}, rev.\ 2021.\newline https://plato.stanford.edu/entries/porphyry/.

\refitem Evangeliou, Christos. ``Aristotle's Doctrine of Predicables and Porphyry's \gri{Isagoge}.'' \gri{Journal of the History of Philosophy} 23 (1985): 15--34.

\refitem Klima, Gyula. ``The Medieval Problem of Universals.'' In \gri{Stanford Encyclopedia of Philosophy}, rev.\ 2022.\newline https://plato.stanford.edu/entries/universals-medieval/.

\refitem Verboon, Annemieke. ``The Medieval Tree of Porphyry.'' In \gri{The Tree: Symbol, Allegory and Structural Device in Medieval Art and Thought}, International Medieval Research 20, 83--101. Turnhout: Brepols, 2014.

\end{references}

\end{document}
